\documentclass{subfiles}
\begin{document}
    Named after its inventors, \citeauthor{RivestSA1977},
        RSA is a public-key schema based on the hardness of integers factorization.
        The core idea is simple: choose two odd prime \(p\) and \(q\) and compute their product \(n\).
        If \(m\) is the message to be transmitted, encode it using \(n\) and send the encoded message.
        \autoref{proc:RSA-keygen} describes how each entity involved computes its keys,
            while \autoref{proc:RSA-encode-decode} shows how a message is encoded/decoded.

        \subfile{../TikZ/Procedure 1 - RSA keygen}
        Here \(e\) is called \emph{encryption key} and \(d\) is called \emph{decryptio key},
        while \(n\) is called \emph{modulo}.

        \subfile{../TikZ/Procedure 2 - RSA encode-decode}
\end{document}
